% ==============================================================================
% Capítulo 8: Considerações finais
% ==============================================================================
\chapter{Considerações finais}
\label{cap:conclusao}

\begin{instrucao}
\orientacao{O que vai nas Considerações Finais}
Retome o problema e os objetivos, mostre como cada objetivo específico foi
atingido e em que capítulo, resuma o que foi entregue, reconheça as
limitações, relate as lições aprendidas e as dificuldades enfrentadas e aponte
trabalhos futuros. Como referência, de duas a quatro páginas costumam ser
suficientes.

Não apresente resultados novos. Em trabalhos tecnológicos, vale informar a
situação do sistema ao final do trabalho (em uso, em homologação, aguardando
implantação) e o que falta para colocá-lo em produção.
\end{instrucao}

Este trabalho apresentou o Monitora, um sistema web para agendamento de
monitorias que substitui o processo informal baseado em mensagens e planilhas.
O sistema permite que monitores publiquem horários, que estudantes reservem
atendimentos e que professores acompanhem a frequência e os assuntos
atendidos, mantendo os dados em infraestrutura da universidade.

Os objetivos específicos foram atingidos da seguinte forma. Os requisitos
foram levantados com monitores e professores e especificados no
\cref{cap:requisitos}. A comparação com as soluções existentes, apresentada no
\cref{cap:solucoes}, mostrou que nenhuma delas atendia ao conjunto de
necessidades identificado. A arquitetura e o modelo de dados foram descritos
no \cref{cap:projeto}. A implementação e a implantação em ambiente de
homologação foram apresentadas no \cref{cap:implementacao}. Por fim, o \cref{cap:testes} mostrou
que todos os requisitos não funcionais foram atendidos e que os usuários
avaliaram positivamente a usabilidade do sistema.

% ------------------------------------------------------------------------------
\section{Limitações e lições aprendidas}
\label{sec:conclusao:limitacoes}

O sistema ainda não está integrado ao sistema acadêmico da universidade, o que
obriga o professor a cadastrar manualmente os monitores a cada semestre. A
avaliação com usuários envolveu um número pequeno de participantes de apenas
duas disciplinas. Por isso, os resultados podem não se repetir em turmas maiores.

Do ponto de vista do processo, a validação dos protótipos antes da
implementação evitou retrabalho, já que duas mudanças importantes na tela de
reserva surgiram nessa etapa. A principal dificuldade foi estimar o esforço das
histórias nos primeiros ciclos, que ficaram sobrecarregados. A partir do
terceiro ciclo, o uso da velocidade observada nos ciclos anteriores tornou o
planejamento mais preciso.

% ------------------------------------------------------------------------------
\section{Trabalhos futuros}
\label{sec:conclusao:futuros}

\begin{instrucao}
\orientacao{Trabalhos futuros}
Evite listas genéricas como ``melhorar a interface''. Para cada sugestão,
explique por que ela é relevante e como poderia ser feita, com base no que o
trabalho mostrou (requisitos de baixa prioridade, comentários dos usuários,
limitações encontradas).
\end{instrucao}

A sugestão mais pedida pelos professores foi exportar os relatórios em
planilha, o que pode ser feito com pouco esforço a partir das consultas já
implementadas. A integração com o sistema acadêmico eliminaria o cadastro
manual de monitores, mas depende da disponibilização de uma interface de
programação pela equipe de tecnologia da universidade. Por fim, um painel com
os assuntos mais procurados em cada disciplina poderia ajudar os professores a
identificar conteúdos que precisam de reforço em sala de aula.
